\section{詳細解説}

\subsection{学習した方策の中身}

軌道生成器型（tg）の方策は、(1) ループの中心・振幅・形と、(2) 姿勢誤差 $e$ から位相の速さ $u$ を決めるフィードバック則を学習した。
学習後のフィードバック則を取り出すと（図~\ref{fig:D6_learned_policy.png} 左）、
\begin{itemize}
\item 誤差が正（目標が反時計回り側）なら $u=+1$（ループを正方向に最大速度で回す）、負なら $u=-1$。
\item 誤差がほぼゼロの狭い範囲で $u$ が急に 0 を横切る（$e=0$ で $u=-0.014$）。
\end{itemize}
つまり学習方策は、\textbf{「目標の向きに応じてループの回し方の向きを切り替え、目標付近で止まる」という、ほぼ最短時間（バンバン）型の制御則}を見つけた。
誤差が $\pm180^\circ$ 付近で $u$ が反転しかけているのは、学習に使った目標（$\pm150^\circ$ まで）の外側であるためで、訓練範囲外では方策が保証されないことを示している。

学習したループの形（図~\ref{fig:D6_learned_policy.png} 右）は、可動域で切り取られた大きな楕円で、曲率 $F<0$ の領域（破線の内側）をほぼ埋めている。

\DFIG{D6_learned_policy.png}{学習した軌道生成器型方策の中身。左：姿勢誤差から位相速度へのフィードバック則。右：学習したループ（緑）と曲率マップ。}

\subsection{学習の評価の全結果}

\DTAB{sim05a_summary}

\subsection{推進剤最小の計画の構造}

スラスタ＋AMBAC の推進剤最小の計画を、旋回時間 20 s・60 s・110 s で比べた（図~\ref{fig:D7_lp_plans.png}）。
\begin{itemize}
\item どの場合も、AMBAC は全区間で最大の周波数 $f_{\max}$ で回す（無料なので使い切る）。
\item スラスタは、最初と最後の 0.5 s 程度だけ噴射する。最初に回転の角運動量を与え、最後に打ち消す（バンバン型）。
\item AMBAC だけで稼げない分の回転を、スラスタの角運動量による等速の回転で補う。旋回時間が長いほど、スラスタの分担は小さくなる。
\end{itemize}
推進剤は 20 s で 9.37 kg、60 s で 1.97 kg、110 s で 0.32 kg である。

\DFIG{D7_lp_plans.png}{推進剤最小の計画（スラスタ＋AMBAC、18 m 級）。上：スラスタのトルク。下：周回平均の姿勢。}

\subsection{パレート曲線の数値}

\DTAB{sim05b_pareto}

\subsection{モデル誤差・外乱下の比較の全結果}

\DTAB{sim05b_eval}

\section*{考察課題}
\begin{exercise}
学習方策のフィードバック則（図~\ref{fig:D6_learned_policy.png} 左）は、誤差がゼロの付近で急に切り替わる。実機で測定ノイズがあると何が起きうるか。対策を考えよ。
\end{exercise}
\begin{exercise}
「計画は最適制御、実行は学習」というハイブリッドを作るとき、学習方策には何を観測させ、何を出力させるのがよいか。考察 C-2（角運動量を一度入れると AMBAC は推進剤を節約しない）を踏まえて提案せよ。
\end{exercise}
